\documentclass[12pt]{extarticle}

\usepackage[paperheight=11in, paperwidth=8.5in, top=2cm, bottom=2cm, left=2cm, right=2cm]{geometry}
\usepackage[thinc]{esdiff}
\usepackage{amsmath}
\usepackage{amssymb}
\usepackage{bm}
%\usepackage{enumitem}
\usepackage{enumerate}
\usepackage{fancyhdr}
\usepackage{lipsum}
\usepackage{hyperref}
\usepackage{textcomp}
\usepackage[toc,page]{appendix}
\usepackage{titlesec}
\usepackage{titletoc}

\newcommand{\mytexttilde}{\raisebox{0.5ex}{\texttildelow}}
\newcommand{\mycomment}[1]{}

\makeatletter
\renewcommand*\env@matrix[1][*\c@MaxMatrixCols c]{%
  \hskip -\arraycolsep
  \let\@ifnextchar\new@ifnextchar
  \array{#1}}
\makeatother


\usepackage{graphicx}
\graphicspath{ {./images/} }

\usepackage{pdfpages}

\hypersetup{
    colorlinks=true,
    linkcolor=blue,
    filecolor=magenta,      
    urlcolor=cyan,
    pdftitle={Lab04},
    pdfpagemode=FullScreen,
    }

\urlstyle{same}
\setlength{\headheight}{15pt}

\begin{document}

\begin{titlepage}
    \begin{center}
        \vspace*{1cm}
            
        \Huge
        \textbf{SFA FECO Analyzer Manual}
            
        \vspace{0.5cm}
        \LARGE
        February 15, 2026
            
        \vspace{1.5cm}
            
        \textbf{v0.0.1}
            
            
    \end{center}
\end{titlepage}	

\tableofcontents

\newpage

\section{Foreword}

The SFA FECO Analyzer is a Python-based tool for analyzing motion in grayscale TIFF images generated by the Surface Fores Apparatus, providing interactive features such as image cropping, wave centerline analysis, and post-analysis data cleaning. The tool allows users to crop the region of interest from the input image, analyze the wave motion across frames, interactively delete unwanted or broken data points from the results, and finally generate force-distance profiles from the data.
The goal of the software is to be easy to use while still allowing for modularity and ability to analyze data in other software. This manual is up to data as of version \textit{2.2.1}.

\subsection{Methods of Installation}

There are three primary methods of installing the software, listed below in order of recommendation:

\begin{enumerate}
	\item Download executable from Github
	\item Pull source code from Github
	\item Compile executable from source code
\end{enumerate}

\subsubsection{Downloading executable}

	Provided on the Github page for this project are multiple compressed files for different operating systems. Currently, executables have been provided for Windows (win86\_64.zip), Mac (mac.zip), and Linux (linux86.tar.gz). If you wish to run the software from a different operating system or architecture, please continue to the section ``\textit{1.1.2 Pulling Source Code}''.
	After downloading and extracting the archive file, there will be a main.exe file and a folder called \_internal. These two files must stay next to each other, and from here, the executable may be executed as normal, without the need to download python or set up any other environment. The only other note is that, should you want to use to software to convert files from CDX files to TIFF files, you must have java already installed on your system.

\section{Usage Information}
\subsection{Dispersion Calculation}

\subsection{Radius of Curvature Calculation}

\subsubsection{Flat Contact Diameter}

\subsection{Preparation Step}

\subsection{Analyze}

\subsection{Visualize}

\subsection{Saving/Loading Calibration Files}

\section{Technical Information}
\subsection{Dispersion Calculation}

\subsection{Radius of Curvature Calculation}

\subsubsection{Flat Contact Diameter}

\subsection{Preparation Step}

\subsection{Analyze}

\subsection{Visualize}

\subsection{Saving/Loading Calibration Files}

\end{document}